\subsection{Núcleo común de las \(3\) formulaciones de dualidad T}

Definamos el dominio estable
\[
 \mathscr D_0=\mathbb Z^2\times\mathbb R_{>0},
 \qquad
 E_{R_0}(m,w)=\frac{m^2}{R_0^2}+w^2R_0^2,
\]
y la involución
\[
 \mathsf T_0(m,w,R_0)=(w,m,R_0^{-1}).
\]

\begin{theorem}[Equivalencia del núcleo ambiente residual--cociente]
\label{thm:cierre-vtm-tres-formulaciones-t}
Después de prolongar las variables \(r,Q\) al dominio estable
\(\mathbb Z^2\), las tres formulaciones se identifican mediante
\[
 I_F(r,Q,R)=(m,w,R_0)=(r,Q,R),
 \qquad I_{16\to02}=\operatorname{id}_{\mathscr D_0}.
\]
Estos mapas entrelazan las involuciones y las formas cuadráticas:
\[
 I_F\mathsf T_F=\mathsf T_0I_F,
 \qquad
 E_R^F=E_{R_0}\circ I_F,
 \qquad
 \mathsf T_0^2=\operatorname{id},
 \qquad
 E_{R_0}=E_{R_0^{-1}}\circ\operatorname{swap}.
\]
\end{theorem}

\begin{proof}
Los dos mapas son renombramientos de las mismas tres coordenadas. Además,
\[
 \mathsf T_0^2(m,w,R_0)=(m,w,R_0)
\]
y
\[
 E_{R_0^{-1}}(w,m)
 =\frac{w^2}{R_0^{-2}}+m^2R_0^{-2}
 =w^2R_0^2+\frac{m^2}{R_0^2}
 =E_{R_0}(m,w).
\]
Las identidades son las mismas tras aplicar \(I_F\) o la identidad entre las
dos formulaciones.
\end{proof}

En \(\mathcal H_{\mathrm{lat}}=\ell^2(\mathbb Z^2)\), definamos
\[
 H(R_0)|m,w\rangle=E_{R_0}(m,w)|m,w\rangle,
 \qquad U_T|m,w\rangle=|w,m\rangle.
\]
Entonces \(U_T\) es unitario, \(U_T^2=I\), y la identidad escalar anterior
equivale, sobre la base ortonormal, a
\[
 \boxed{U_TH(R_0)U_T^{-1}=H(R_0^{-1}).}
\]
Ésta es la segunda realización exacta del mismo núcleo, no una nueva
hipótesis física.

\subsection{Distinción de las 2 secciones residuales por sus inserciones energéticas y clases de equivalencia}

La convención residual positiva usa
\[
 r_9^+(n)=1+((n-1)\bmod9),
 \qquad Q_9^+(n)=\frac{n-r_9^+(n)}9,
\]
mientras que la convención residual cero inserta la clase
\(\bar r_9(n)\in\{0,\ldots,8\}\). El cambio de sección que conserva el
entero es
\[
 C(r,Q)=\bigl(\bar r,Q+\mathbf1_{\{r=9\}}\bigr).
\]

\begin{proposition}[Obstrucción energética del cambio de sección]
\label{prop:cierre-vtm-obstruccion-secciones-nonadicas}
El mapa \(C\) conserva \(n=9Q+r\), pero no conserva en general la forma
\(E_R\) ni entrelaza la involución de intercambio. Por tanto, la equivalencia
del núcleo ambiente no implica la equivalencia de las dos inserciones
nonádicas.
\end{proposition}

\begin{proof}
Para \(n=9\), la primera convención da \((r,Q)=(9,0)\), mientras que el cambio con acarreo da
\((\bar r,Q_0)=(0,1)\). Sus energías son
\[
 E_R(9,0)=\frac{81}{R^2},
 \qquad E_R(0,1)=R^2,
\]
que no coinciden para todo \(R>0\). La inserción sin corregir
el acarreo enviaría este ejemplo a \((0,0)\) y tampoco conservaría el entero.
Se necesita fijar una única sección o construir una forma covariante bajo el
cambio de sección antes de promover esta equivalencia.
\end{proof}

\begin{definition}[Cociente de cambio de sección y energía covariante]
\label{def:cierre-vtm-energia-covariante-seccion}
Sea
\[
 L_9:=\mathbb Z(9,-1)\subset\mathbb Z^2,
 \qquad
 S(x_1,x_2)=(x_2,x_1),
\]
y sea
\[
 q_R(x_1,x_2)=\frac{x_1^2}{R^2}+x_2^2R^2.
\]
Sobre los cocientes
\[
 \mathscr Q_9=(\mathbb Z^2/L_9)\times\mathbb R_{>0},
 \qquad
 \mathscr Q_9^\vee=(\mathbb Z^2/SL_9)\times\mathbb R_{>0},
\]
definimos
\[
 \mathcal E_R^{L_9}([x]_{L_9})
 :=\min_{\ell\in L_9}q_R(x+\ell)
\]
y
\[
 \mathsf T_{\mathrm{cov}}([x]_{L_9},R)
 :=([Sx]_{SL_9},R^{-1}).
\]
\end{definition}

\begin{theorem}[Equivalencia covariante de las secciones nonádicas]
\label{thm:cierre-vtm-equivalencia-covariante-secciones}
La energía \(\mathcal E_R^{L_9}\) y la transformación
\(\mathsf T_{\mathrm{cov}}\) están bien definidas. Las convenciones
\((9,0)\) y \((0,1)\), y en general las secciones \(1,\ldots,9\) y
\(0,\ldots,8\) con el acarreo correspondiente, representan la misma clase
de \(\mathbb Z^2/L_9\). Además,
\[
 \boxed{
 \mathcal E_{R^{-1}}^{SL_9}([Sx]_{SL_9})
 =\mathcal E_R^{L_9}([x]_{L_9}).}
\]
El intercambio inverso, definido sobre \(\mathscr Q_9^\vee\), devuelve la
clase inicial; por tanto las dos cartas cocientadas son T-duales.
\end{theorem}

\begin{proof}
Si \(x'=x+\ell_0\), la traslación
\(\ell\mapsto\ell+\ell_0\) permuta \(L_9\), luego el mínimo no depende del
representante. Existe porque \(q_R(x+k(9,-1))\) es una función cuadrática
coerciva de \(k\in\mathbb Z\). La diferencia
\((9,0)-(0,1)=(9,-1)\) pertenece a \(L_9\), y lo mismo ocurre en cada cambio
de sección con acarreo. Finalmente,
\[
\begin{aligned}
 \mathcal E_{R^{-1}}^{SL_9}([Sx])
 &=\min_{\ell\in L_9}q_{R^{-1}}(Sx+S\ell)\\
 &=\min_{\ell\in L_9}q_R(x+\ell)
 =\mathcal E_R^{L_9}([x]).
\end{aligned}
\]
Como \(S^2=I\), el segundo intercambio recupera el cociente original y el
radio \(R\).
\end{proof}

\paragraph{Dominio de equivalencia del núcleo dual y no equivalencia de las inserciones energéticas.}
La equivalencia de los núcleos sobre \(\mathscr D_0\) es exacta. Las dos
inserciones no son equivalentes para la energía no corregida; el cierre
covariante es el cociente y la energía del teorema anterior.

\subsection{Normalización exacta de la realización de cuerda}

Fijemos \(\alpha'>0\). En el dominio
\(\mathscr D_{\alpha'}=\mathbb Z^2\times\mathbb R_{>0}\), sea
\[
 \mathsf T_{\alpha'}(m,w,R)=\left(w,m,\frac{\alpha'}R\right),
 \qquad
 \mathcal N_{\alpha'}(m,w,R)
 =\left(m,w,\frac R{\sqrt{\alpha'}}\right).
\]

\begin{theorem}[Conjugación por la escala de cuerda]
\label{thm:cierre-vtm-normalizacion-alpha-t}
La normalización \(\mathcal N_{\alpha'}\) es una conjugación exacta:
\[
 \boxed{\mathcal N_{\alpha'}\mathsf T_{\alpha'}
 =\mathsf T_0\mathcal N_{\alpha'}.}
\]
Para
\[
 E_{\mathrm{str}}(m,w,R)
 =\frac{m^2}{R^2}+\frac{w^2R^2}{\alpha'^2},
\]
se cumple
\[
 E_{\mathrm{str}}
 =\alpha'^{-1}E_0\circ\mathcal N_{\alpha'}.
\]
Además, si
\[
 p_L=\frac mR+\frac{wR}{\alpha'},
 \qquad p_R=\frac mR-\frac{wR}{\alpha'},
\]
entonces \(p_L\) queda fijo y \(p_R\mapsto-p_R\).
\end{theorem}

\begin{proof}
Se tiene
\[
 \mathcal N_{\alpha'}\mathsf T_{\alpha'}(m,w,R)
 =\left(w,m,\frac{\sqrt{\alpha'}}R\right)
 =\mathsf T_0\mathcal N_{\alpha'}(m,w,R).
\]
La identidad de energías resulta de
\[
 E_0\!\left(m,w,\frac R{\sqrt{\alpha'}}\right)
 =\frac{\alpha'm^2}{R^2}+\frac{w^2R^2}{\alpha'}
 =\alpha' E_{\mathrm{str}}(m,w,R).
\]
La sustitución de
\((m,w,R)\mapsto(w,m,\alpha'/R)\) en \((p_L,p_R)\) da, respectivamente,
\(p_L\) y \(-p_R\).
\end{proof}

La conjugación y la invariancia del espectro con datos oscilatorios fijos son
exactas. La identificación de \(m,w,R,\alpha'\) con momento, enrollamiento,
radio físico y tensión inversa es una interfaz física con premisas propias.

\subsection{Nivel, intercambio lineal de masas y modularidad: \(3\) controles de no fusión}

La condición activa de cuerda es
\[
 N_L-N_R=a_L-a_R-mw.
\]
La fórmula esquemática \(N_L-N_R=nw\) coincide con ella sólo
después de declarar \(a_L=a_R\) y \(n=-m\), o una inversión equivalente de
la orientación de hojas. Sin ese diccionario, las dos expresiones no son la
misma formulación.

La simetría lineal de la acción se tipa en el espacio ampliado de
coeficientes y modos:
\[
 L(A,C;n,w)=nA+wC,
 \qquad
 \sigma(A,C;n,w)=(C,A;w,n).
\]
Entonces
\[
 L\circ\sigma=wC+nA=L.
\]
Es una identidad exacta de intercambio simultáneo de parámetros y
coordenadas. No equivale a la invariancia cuadrática de \(H(R)\) sobre un
sector con parámetros físicos fijos; tal promoción requiere un
entrelazador adicional.

Finalmente, sobre el eje imaginario del semiplano superior, definamos
\[
 \iota(R_0)=iR_0^2,
 \qquad S(\tau)=-\frac1\tau.
\]
Entonces
\[
 \boxed{S\circ\iota(R_0)=\iota(R_0^{-1}).}
\]
Como \(J(S\tau)=J(\tau)\), se obtiene
\(J(iR_0^2)=J(iR_0^{-2})\). Ésta es una semiconjugación exacta del parámetro
de radio con la transformación modular \(S\), demostrada con la estructura
modular de partida explícita. No identifica por sí sola la dualidad de
espacio objetivo, el intercambio de cargas y la construcción de
\(V^\natural\).
